
\subsection{RQ8: What are \toolname's dominant failure modes?}
\label{subsec:rq8-failure}

Fig.~\ref{fig:failure_root_causes} analyzes an independent
\texttt{gpt-5.4-mini} rerun with per-case logs. Capping refinement at three
iterations instead of ten exposes more failures: 321 valid cases versus 373
in the main run. The panels show terminal outcomes and the dominant program
feature in each of the 79 non-valid cases (42 \emph{No valid}, 37 \emph{No
syntax}); for overlapping features, we choose the first explaining the
blocking obligation. Table~\supref{tab:gpt54mini_failure_events}
(\appref{sec:limitations}) separately classifies logged syntax and prover
events by outcome and obligation type.

Loop invariant obligations account for \(2{,}590/4{,}370=59.3\%\) of
timeouts, especially establishment and preservation over arrays and heap
arrays. Postcondition and \texttt{assigns} timeouts also occur.

\rev{Log outcomes \emph{Timeout}, \emph{Unknown}, and \emph{Failed} denote
timeouts, prover uncertainty, and backend or resource failures, not semantic
counterexamples. They cannot separate incorrect specifications from correct
but unproved obligations. Manual review more often found plausible properties
limited by complex loop, array, or heap proofs, without ruling out synthesis
errors.} See \rev{\appref{sec:limitations}} for method and verifier limits.

\begin{figure}[t]
\centering
\begin{minipage}[t]{\columnwidth}
\centering
\begin{tikzpicture}[font=\footnotesize,every node/.style={text=figgraytext}]
  \path[use as bounding box] (-1.20,-1.10) rectangle (6.25,1.10);
  \def\r{1.00}
  \filldraw[draw=figgrayframe,line width=0.3pt,fill=figgrayfill]
    (0,0) -- (0:\r) arc[start angle=0,end angle=288.9,radius=\r] -- cycle;
  \filldraw[draw=figcontractstroke,line width=0.3pt,fill=figcontractfill]
    (0,0) -- (288.9:\r) arc[start angle=288.9,end angle=326.7,radius=\r] -- cycle;
  \filldraw[draw=figcallstroke,line width=0.3pt,fill=figcallfill]
    (0,0) -- (326.7:\r) arc[start angle=326.7,end angle=360,radius=\r] -- cycle;
  \draw[figgrayframe,line width=0.3pt] (0,0) circle (\r);
  \filldraw[fill=figgrayfill,draw=figgrayframe,line width=0.25pt] (1.55,0.54) rectangle +(0.16,0.16);
  \node[anchor=west] at (1.82,0.62) {Valid: 321 (80.3\%)};
  \filldraw[fill=figcontractfill,draw=figcontractstroke,line width=0.25pt] (1.55,0.22) rectangle +(0.16,0.16);
  \node[anchor=west] at (1.82,0.30) {No valid: 42 (10.5\%)};
  \filldraw[fill=figcallfill,draw=figcallstroke,line width=0.25pt] (1.55,-0.10) rectangle +(0.16,0.16);
  \node[anchor=west] at (1.82,-0.02) {No syntax: 37 (9.2\%)};
\end{tikzpicture}
\par\smallskip{\footnotesize\normalfont (a) Terminal outcomes.}
\end{minipage}
\par\smallskip
\begin{minipage}[t]{\columnwidth}
\centering
\begin{tikzpicture}[font=\footnotesize,every node/.style={text=figgraytext}]
  \path[use as bounding box] (-1.20,-1.15) rectangle (6.25,1.15);
  \def\r{1.00}
  \filldraw[draw=figgrayframe,line width=0.28pt,fill=figgrayfill]
    (0,0) -- (0:\r) arc[start angle=0,end angle=141.27,radius=\r] -- cycle;
  \filldraw[draw=figloopstroke,line width=0.28pt,fill=figloopfill]
    (0,0) -- (141.27:\r) arc[start angle=141.27,end angle=250.63,radius=\r] -- cycle;
  \filldraw[draw=figgrayframe,line width=0.28pt,fill=figgray!10]
    (0,0) -- (250.63:\r) arc[start angle=250.63,end angle=314.43,radius=\r] -- cycle;
  \filldraw[draw=figcallstroke,line width=0.28pt,fill=figcallfill]
    (0,0) -- (314.43:\r) arc[start angle=314.43,end angle=341.77,radius=\r] -- cycle;
  \filldraw[draw=figloopstroke,line width=0.28pt,fill=figloop!20]
    (0,0) -- (341.77:\r) arc[start angle=341.77,end angle=350.89,radius=\r] -- cycle;
  \filldraw[draw=figgrayframe,line width=0.28pt,fill=figgray!6]
    (0,0) -- (350.89:\r) arc[start angle=350.89,end angle=355.44,radius=\r] -- cycle;
  \filldraw[draw=figgrayframe,line width=0.28pt,fill=white]
    (0,0) -- (355.44:\r) arc[start angle=355.44,end angle=360,radius=\r] -- cycle;
  \draw[figgrayframe,line width=0.3pt] (0,0) circle (\r);
  \filldraw[fill=figgrayfill,draw=figgrayframe,line width=0.2pt] (1.45,0.895) rectangle +(0.13,0.13);
  \node[anchor=west] at (1.68,0.96) {Heap dyn. arrays: 31 (39.2\%)};
  \filldraw[fill=figloopfill,draw=figloopstroke,line width=0.2pt] (1.45,0.575) rectangle +(0.13,0.13);
  \node[anchor=west] at (1.68,0.64) {Array pointer loops: 24 (30.4\%)};
  \filldraw[fill=figgray!10,draw=figgrayframe,line width=0.2pt] (1.45,0.255) rectangle +(0.13,0.13);
  \node[anchor=west] at (1.68,0.32) {String char logic: 14 (17.7\%)};
  \filldraw[fill=figcallfill,draw=figcallstroke,line width=0.2pt] (1.45,-0.065) rectangle +(0.13,0.13);
  \node[anchor=west] at (1.68,0.00) {Recursive callee assert: 6 (7.6\%)};
  \filldraw[fill=figloop!20,draw=figloopstroke,line width=0.2pt] (1.45,-0.385) rectangle +(0.13,0.13);
  \node[anchor=west] at (1.68,-0.32) {Scalar loops: 2 (2.5\%)};
  \filldraw[fill=figgray!6,draw=figgrayframe,line width=0.2pt] (1.45,-0.705) rectangle +(0.13,0.13);
  \node[anchor=west] at (1.68,-0.64) {Floating-point casts: 1 (1.3\%)};
  \filldraw[fill=white,draw=figgrayframe,line width=0.2pt] (1.45,-1.025) rectangle +(0.13,0.13);
  \node[anchor=west] at (1.68,-0.96) {Other: 1 (1.3\%)};
\end{tikzpicture}
\par\smallskip{\footnotesize\normalfont (b) Root causes for non-valid cases.}
\end{minipage}
\caption{Failure outcome and root-cause summary for the \texttt{gpt-5.4-mini} \toolname-400 rerun.}
\label{fig:failure_root_causes}
\end{figure}

\begin{finding}
\textbf{Finding (RQ8).} \rev{Residual failures manifest predominantly as SMT timeouts on complex proof obligations. The logs cannot conclusively separate specification-synthesis errors from proof obligations beyond the current backend, although manual inspection more often points to verifier limitations.}
\end{finding}
